\chapter{Preface}
\label{chap:preface}

The purpose of this thesis is to produce a tutorial. A working system was built in the course of it, and that system is described here in full, though the intended result of the work is not only the system itself but also the documented procedure that leads to it. The tutorial form is therefore an objective rather than a stylistic decision, and it governs how the material is arranged. Every step is given in the order in which it must be performed, so that a reader starting from a bare board and an empty toolchain can reproduce the whole result.

The system to be built is a complete video system on a single \gterm{socfpga}{SoC FPGA}, with acquisition from a camera carried out in hardware, and processing carried out in software, on the processor system integrated in the same device.

\section{The hardware provided}

The assigned platform is a \term{Terasic DE1-SoC} development board, built around an \term{Intel Cyclone V SoC}. This device is an \gterm{fpga}{FPGA} that integrates a processor: on the same \gterm{die}{die} sit the reconfigurable logic and a hard processor system (\gterm{hps}{HPS}), a dual-core \term{ARM Cortex-A9} with its own \term{DDR3} memory controller, joined by dedicated \term{HPS-to-FPGA} bridges. Custom logic can therefore be addressed by software as ordinary memory-mapped registers and can in turn reach the processor's own \term{DDR3} memory directly. Almost every design decision described in the following chapters follows from this arrangement.

The image sensor is an \gterm{ov7670}{OV7670} camera module, a 640$\times$480 resolution \term{CMOS} rolling-shutter device capable of 30 frames per second and able to code each pixel as \gterm{bayermosaic}{Bayer mosaic}, \gterm{rgb565}{RGB565}, \gterm{rgb555}{RGB555}, \gterm{rgb444}{RGB444} or \gterm{yuv422}{YUV 4:2:2}. Its transmission protocol is strictly unidirectional, with no handshaking and no flow control, so pixels leave the sensor on an 8-bit parallel bus timed by a pixel clock and delimited by line and frame signals. The receiver then has to follow the sensor's timing rather than ask it for data. The pixel clock is not free-running, since it is derived from a \uline{master clock} which must be supplied from outside. The sensor also exposes an \gterm{i2c}{I\textsuperscript{2}C}-like control bus, which is used to set up resolution, output format and the other imaging parameters at run time by writing its configuration registers. Reading the video stream from such a device is therefore not a trivial connection to a well-known interface, but the construction of a custom one whose logic interprets its timing and forwards the data to the \gterm{hps}{HPS} through the \gterm{fpga}{FPGA} fabric.

During development it became clear that the supplied jumper leads, about ten centimetres long, degraded the video signal, and that the twenty-centimetre ones, needed to move the camera freely, corrupted it altogether. For this reason, a pair of small \gterm{pcb}{PCB}s, one at each end of the cable, was designed and built to drive the data line. Generating the camera \uline{master clock} in the \gterm{fpga}{FPGA} was weighed at that point and set aside, since it would have called for an additional line buffer, and a cheaper clock generator sitting next to the camera served instead.

\section{Requirements and constraints}

Two requirements shaped the development. The first, set by the project, is to relieve the processor of the image data transfer, or at least to involve it as little as possible. Its computing resources are meant for computing, not for carrying frames from the acquisition logic into memory: keeping a processor busy with a sustained transfer that involves no computation at all wastes the \term{CPU}'s cycles and energy alike.

The second, self-imposed and softer, is that no step of the procedure depends on a paid licence. Every tool used here, from the \gterm{fpga}{FPGA} design software to the compiler and the host utilities, is free of charge or open source, so that anyone taking this work further can go from a bare board to a running system without buying anything.

\section{Outline of the implementation}

\subsection{The hardware}

Given the first requirement, the natural answer is a \gterm{dma}{DMA} unit, able to carry the data on its own, straight from the acquisition logic into the processor's \term{DDR3} memory. Thus, all the software has to supply is the address at which the incoming frame is to be stored.

The \gterm{cyclonevsoc}{Cyclone V SoC} at the centre of the board provided offers no such unit ready in the \gterm{hps}{HPS}. The \gterm{dma}{DMA} must therefore be implemented in the \gterm{fpga}{FPGA} fabric and attached to the \gterm{hps}{HPS}'s \term{DDR3} controller through the \term{HPS-to-FPGA} bridges it exposes.

This can be accomplished with the manufacturer's own design software, \gterm{quartusprime}{Quartus Prime}, which includes \gterm{platformdesigner}{Platform Designer}, a tool for composing a system out of hardware blocks.\footnote{\gterm{platformdesigner}{Platform Designer} treats the \gterm{hps}{HPS} itself as one such block.} They may be drawn from the vendor's own library or written for the occasion: each one must only present a compatible interface. \gterm{platformdesigner}{Platform Designer} then generates the interconnect that joins them together, in the form of hardware description language. \term{Intel} supplies a \term{Modular Scatter-Gather DMA}, or \gterm{msgdma}{mSGDMA}, which fits the requirement exactly and is the unit used throughout this work. The remaining modules that complete the capture path are realized as \gterm{platformdesigner}{Platform Designer} components too, custom ones when the library held nothing that did the job.

The generic part of the hardware, essentially everything outside the video path, follows two reference designs, a published design for the \term{DE1-SoC} \cite{UPM-DE1} and \term{Intel}'s \term{Golden Hardware Reference Design} \cite{GHRD}, the latter targeting a different board built around a similar \gterm{cyclonevsoc}{Cyclone V SoC}. Neither could be reused as it stood, one because it targets another board and the other because it has no video path, but together they cover the obligatory groundwork, how the system is instantiated and clocked, how reset is distributed and how the pins of the board are assigned. That part of the design admits one correct solution, so there was nothing to invent there.

\subsection{The software}

The software side\footnote{And the related hardware, since each repository carries the \term{Quartus} project alongside the software.} took the form of two projects, kept in separate repositories, \gterm{fpgalix}{\prjname{FPGAlix}}\cite{FPGALIX-REPO}, which runs under embedded \term{Linux} on the \gterm{hps}{HPS}, and \prjname{FPGAsteel}\cite{FPGASTEEL-REPO}, which runs \gterm{baremetal}{bare metal} on the same processor. The two environments differ in what they make affordable. For bringing the board up, \term{Linux} offers a safe road, since the \gterm{bootchain}{boot chain}, the drivers and the capture framework are code that has been made to work before, so when the system does not come up the fault lies in how it has been configured rather than in the software itself. For the application it supplies device drivers, memory management, \gterm{cache}{cache} management, a standard video capture framework and tested libraries, which shorten the work, and it allows the program to be debugged remotely over \term{Ethernet} with \cmdpath{gdb}, while bare metal has to be followed through \term{USB} \gterm{jtag}{JTAG}. What it costs is distance from the hardware and timing that is not entirely predictable. Bare metal removes the abstraction and the overhead and gives complete control of the entire system, but everything must be written from nothing, and the problems the operating system hides come back into view. This extends to the processing code itself: routines that under \term{Linux} could draw on existing library functions must here be implemented from first principles.

From the hardware point of view the two projects differ slightly. In both a single frame can be read out through the debugger and examined on the development host, while the display of a continuous video stream requires a dedicated output path, and the two projects adopt different ones. \prjname{FPGAlix} transmits the frames over \term{Ethernet}, relying on the interface and the network stack provided by the operating system, at the cost of a significant processor load, due to encoding and to the splitting of the stream into network packets. \prjname{FPGAsteel} drives the \gterm{vga}{VGA} port available on the board,\footnote{This path too is built around an \gterm{msgdma}{mSGDMA}, which reads the frame back from memory and streams it to the video output.} since implementing a network driver and a \term{TCP/IP} stack in bare metal would have required an effort disproportionate to the aim of the project.

\section{Organization of the document}

The two projects just introduced, \prjname{FPGAlix} and \prjname{FPGAsteel}, are described in the two parts that follow in the tutorial form the aim of the work requires, with every step given in the order in which it must be carried out. A working familiarity with \gterm{quartusprime}{Quartus Prime} and its \gterm{platformdesigner}{Platform Designer}, with \gterm{cpp}{C++} and with the \term{Linux} command line is assumed. A reader who lacks the first two will find in \cite{UPM-DE1} a guided introduction which, although written for an older release of the tools, still describes the \term{Quartus} side closely enough to be followed.

\textbf{Part~\ref{part:fpgalix}} covers \prjname{FPGAlix}. It defines the first version of the hardware, brings up the operating system on the \gterm{hps}{HPS} and closes with one example of video processing. A single example is enough here, because under the operating system the software is divided into modules with well-defined boundaries, each of them behaving in a familiar and predictable way, so further examples would add length without adding understanding.

\textbf{Part~\ref{part:fpgasteel}} covers \prjname{FPGAsteel}. On the hardware side it adds the \gterm{dma}{DMA} that feeds the \gterm{vga}{VGA} video output, and on the software side it offers several examples. Bare-metal code is tailored to a single system and very little of it is standard since it heavily depends on the choices of whoever writes it. The examples therefore proceed by steps, each one adding a single capability to the one before it:

\begin{enumerate}
    \item prints a message on the \gterm{serialconsole}{serial console},
    \item establishes an \gterm{i2c}{I\textsuperscript{2}C} communication for configuring the camera,
    \item forwards the incoming video stream from the camera to \gterm{vga}{VGA} output,
    \item tests the same path with the caches enabled,
    \item brings up the second core,
    \item implements an image-processing chain, which runs on both cores, for locating the boundaries and the outline of a bright object against a dark background.
\end{enumerate}

Part~\ref{part:appendices} gathers the appendices: the schematics and the bill of materials of the camera adapter board, the glossary of terms and acronyms used throughout the document, and the bibliography, listing the project repositories, the external documents referred to throughout and the artefacts produced by the work.

The first two parts, however, can also be read separately, by a reader interested only in the \term{Linux} side or only in the bare-metal one.

\section{Source repositories and inline-code conventions}

Two public repositories accompany this thesis, \prjname{FPGAlix}\cite{FPGALIX-REPO} and \prjname{FPGAsteel}\cite{FPGASTEEL-REPO}. They hold the implementations described in the chapters that follow: the \term{Quartus} projects, the \gterm{vhdl}{VHDL} sources, the \gterm{cpp}{C++} software and the \gterm{make}{\cmdpath{make}} build scripts. What the document adds is the reasoning behind them, the role each component plays and the order in which the whole is assembled. Unless otherwise stated, a repository-relative path belongs to the project discussed in the current part.

Because the document and the sources are meant to be read together, inline technical identifiers are styled according to what they denote, rather than merely as code. This distinguishes, for example, a repository path from one outside the repositories, a program built here from a generic command, and a hardware register from a source-code identifier. Table~\ref{tab:preface-inline-styles} lists the eight styles used for code-styled text, and the three markers used for an ordinary technical term, for one defined in the Glossary (Appendix~\ref{app:glossary}), and for the name of a project.

\begin{table}[H]
\centering
\begin{tabularx}{\textwidth}{@{} l X @{}}
\toprule
\textbf{Example} & \textbf{What it marks} \\
\midrule
\extpath{/dev/ttyUSB0} & a filesystem path outside the two repositories \\
\repopath{sw/} & a filesystem path inside one of them \\
\lightmidrule
\cmdpath{make} & an executable of the system or of the toolchain \\
\cmd{vision_core} & an executable, script or example built from one of them \\
\lightmidrule
\ccode{Core1::dispatch()} & \term{C/C++} source: functions, classes, members and project-level macros, including register names where they appear inside a code snippet \\
\lightmidrule
\hw{RSTMGR_MPUMODRST} & peripherals, registers and \gterm{platformdesigner}{Platform Designer} signals, where they are named on their own rather than inside a code snippet \\
\lightmidrule
\sysid{PATH} & environment variables, users and groups \\
\lightmidrule
{\bfseries\ttfamily\small python3-numpy} & anything else set in code style without a more specific category, such as \gterm{apt}{apt} or library package names \\
\midrule
\term{DDR3} & a technical term whose definition is not in the Glossary \\
\gterm{jtag}{JTAG} & a term defined in the Glossary (Appendix~\ref{app:glossary}); underlined, and footnoted with that definition at its first occurrence in each part \\
\lightmidrule
\prjname{FPGAlix} & one of the two projects documented here, and of the repositories that hold them \\
\bottomrule
\end{tabularx}
\caption{Inline styles used for code-styled text and technical terms.}
\label{tab:preface-inline-styles}
\end{table}
